\section{Method}
\label{sec:method}

By Section~\ref{sec:prelim} it suffices to bound the volume of every
$P \in \mathcal{P}_{4}(7)$ of type $C_{4}(7)$, and of every
$P \in \mathcal{P}_{4}(8)$ of each simplicial type other than the
cross-polytope. For one type at a time, we write $d!$ times the volume
as a polynomial $D$ in the coordinates of the vertices, bound $D$ on
$(\mathbb{S}^{d-1})^{n}$ by an algebraic identity whose right-hand side is
visibly non-negative there, and check that identity in exact arithmetic.

\subsection{The signed volume}

\begin{definition}
  \label{def:signed-volume}
  An \emph{orientation} of a finite set is an ordering of its elements up to
  even permutations. Let $\mathcal{C} = (n, \mathcal{F})$, where
  $\mathcal{F}$ is a finite set of ordered $d$-tuples
  $F = (F_{1}, \ldots, F_{d})$ of distinct elements of $\{1, \ldots, n\}$,
  called facets. Deleting $F_{a}$ gives a ridge,
  with the orientation of $(F_{1}, \ldots, \widehat{F_{a}}, \ldots, F_{d})$
  if $a$ is even and the opposite one if $a$ is odd. The orientation of
  $\mathcal{C}$ is \emph{coherent} if every ridge lies in exactly two facets
  and these induce opposite orientations on it. A point
  $p = (p_{1}, \ldots, p_{n}) \in (\mathbb{R}^{d})^{n}$ \emph{realises}
  $\mathcal{C}$ if $P = \conv\{p_{1}, \ldots, p_{n}\}$ has the $p_{i}$ as its
  vertices and its facets are exactly the simplices
  $\conv\{p_{j} : j \in F\}$, $F \in \mathcal{F}$. The \emph{signed volume} of
  $p$ is
  \begin{equation}
    \label{eq:signed-volume}
    D(p) \;=\; \sum_{F \in \mathcal{F}}
      \det\bigl(p_{F_{1}}, \ldots, p_{F_{d}}\bigr).
  \end{equation}
\end{definition}

Writing $x_{i,a}$ for the $a$-th coordinate of $p_{i}$ and
$x_{i} = (x_{i,1}, \ldots, x_{i,d})$, $D$ is a polynomial of degree $d$ in the
$dn$ variables $x_{i,a}$ with integer coefficients \artref{signedvolume}.

\begin{lemma}
  \label{lem:volume}
  If the orientation of $\mathcal{C}$ is coherent and $p$ realises
  $\mathcal{C}$, then $d!\,\vol_{d}(P) = \lvert D(p)\rvert$.
\end{lemma}

\begin{proof}
  Fix $y \in \mathbb{R}^{d}$ and expand
  $\det(p_{F_{1}} - y, \ldots, p_{F_{d}} - y)$ by multilinearity. Terms
  containing $y$ twice vanish. A term containing $y$ once is, up to sign, the
  determinant of $y$ and the vertices of a ridge of $F$ in their induced
  order, and by coherence the two facets at that ridge contribute it with
  opposite signs. Hence
  $D(p) = \sum_{F} \det(p_{F_{1}} - y, \ldots, p_{F_{d}} - y)$ for every $y$.
  Take $y$ in the interior of $P$. The simplices $\conv(\{y\} \cup F)$ tile
  $P$, and each summand is $\pm d!$ times the volume of one of them. The sign
  is $+$ exactly when the ordering of $F$ agrees with the orientation that $P$
  induces on its boundary. That orientation is coherent. Any two facets of
  $P$ are joined by a sequence of facets in which consecutive ones share a
  ridge, and along such a sequence a coherent orientation is determined by
  its value on the first facet. Hence the orientation of $\mathcal{C}$ agrees
  with that of the boundary on every facet or on none, and the signs are all
  $+$ or all $-$.
\end{proof}

\begin{remark}[{\cite{berman-hanes-1970}, \cite[Lemma~3.5]{hoehner-ledford-2026}}]
  For $d = 3$, Berman and Hanes state this identity without proof, and Hoehner
  and Ledford prove it with the divergence theorem.
\end{remark}

\subsection{Certificates}

To prove $D \leq \beta$ on $(\mathbb{S}^{d-1})^{n}$, we write $\beta - D$ as a
polynomial that is non-negative on $\mathbb{R}^{dn}$ plus a polynomial that
vanishes on $(\mathbb{S}^{d-1})^{n}$.

\emph{Sums of squares.} A \emph{monomial} is a product of the variables
$x_{i,a}$, repetitions allowed; the empty product $1$ is the monomial of
degree $0$. Let $v(x)$ be the column vector of all monomials of degree at
most $2$, in a fixed order,
\[
  v(x) = \bigl(1,\; x_{1,1}, \ldots, x_{n,d},\; x_{1,1}^{2},\;
  x_{1,1}x_{1,2}, \ldots, x_{n,d}^{2}\bigr)^{\mathsf{T}},
\]
which has $1 + dn + \binom{dn+1}{2}$ entries.

Every polynomial $g$ of degree at most $2$ is $c^{\mathsf{T}} v(x)$ for a
unique vector $c$, and then $g^{2} = v^{\mathsf{T}} (cc^{\mathsf{T}})\, v$.
Hence the sums of squares of such polynomials are exactly the polynomials
\[
  v^{\mathsf{T}} Q\, v = \lvert L^{\mathsf{T}} v\rvert^{2},
  \qquad Q = LL^{\mathsf{T}} \text{ symmetric positive semidefinite}.
\]
Products of two entries of $v$ have degree at most $4$, the degree of $D$.

\emph{Multiples of the sphere equations.} A sum of squares alone cannot
equal $\beta - D$. Since $D(tx) = t^{d} D(x)$, the polynomial $\beta - D$
takes negative values on $\mathbb{R}^{dn}$. The bound holds only because
\[
  \lvert x_{i}\rvert^{2} = x_{i,1}^{2} + \cdots + x_{i,d}^{2} = 1,
  \qquad i = 1, \ldots, n,
\]
so the identity has to use these equations.

For any polynomials $h_{1}, \ldots, h_{n}$, the polynomial
\[
  \sum_{i=1}^{n} h_{i}(x)\bigl(\lvert x_{i}\rvert^{2} - 1\bigr)
\]
vanishes on $(\mathbb{S}^{d-1})^{n}$.

\begin{definition}
  \label{def:certificate}
  A \emph{certificate} for $\mathcal{C}$ and $\beta \in \mathbb{Q}$ is a
  symmetric positive definite matrix $Q$ with rational entries together with
  polynomials $h_{1}, \ldots, h_{n}$ with rational coefficients such that
  \begin{equation}
    \label{eq:certificate}
    \beta - D(x) \;=\; v(x)^{\mathsf{T}} Q\, v(x)
    \;+\; \sum_{i=1}^{n} h_{i}(x)\bigl(\lvert x_{i}\rvert^{2} - 1\bigr)
  \end{equation}
  holds identically as polynomials in the $x_{i,a}$ \artref{certificate}.
\end{definition}

Because \eqref{eq:certificate} is an identity of polynomials with rational
data, it can be verified in exact arithmetic by comparing coefficients
(Section~\ref{sec:verification}).

\begin{proposition}
  \label{prop:certificate}
  Suppose the orientation of $\mathcal{C}$ is coherent and a certificate for
  $\mathcal{C}$ and $\beta$ exists. Then $\lvert D(p)\rvert \leq \beta$ for
  every $p \in (\mathbb{S}^{d-1})^{n}$, and every polytope inscribed in
  $\mathbb{S}^{d-1}$ that realises $\mathcal{C}$ has volume at most
  $\beta/d!$.
\end{proposition}

\begin{proof}
  Let $p \in (\mathbb{S}^{d-1})^{n}$. In \eqref{eq:certificate} the second sum
  vanishes at $p$ and the first is non-negative, so $D(p) \leq \beta$. Let
  $\rho$ be a reflection of $\mathbb{R}^{d}$. Then
  $\rho p = (\rho p_{1}, \ldots, \rho p_{n})$ also lies in
  $(\mathbb{S}^{d-1})^{n}$, and $D(\rho p) = \det(\rho)\,D(p) = -D(p)$, so
  $-D(p) \leq \beta$ as well. Hence $\lvert D(p)\rvert \leq \beta$, and
  Lemma~\ref{lem:volume} gives the volume bound.
\end{proof}

\begin{remark}[{\cite{parrilo-2003}, \cite{putinar-1993}}]
  Identities of the form \eqref{eq:certificate} are the sum-of-squares
  certificates of Parrilo. Their existence is not automatic. If $\beta$ is
  below the maximum of $D$ on $(\mathbb{S}^{d-1})^{n}$ there is none, by
  Proposition~\ref{prop:certificate}. If $\beta$ is above it, Putinar's
  theorem gives an identity of this form with $Q$ positive semidefinite, but
  with monomials of some degree that it does not bound. We use monomials of
  degree at most $2$, the smallest degree for which $v^{\mathsf{T}} Q\, v$
  reaches the degree $4$ of $D$, and nothing guarantees a certificate of this
  degree in advance: the ones used here exist because they were found and
  checked. Positive semidefiniteness of $Q$ would suffice, but we require
  positive definiteness because, by Sylvester's criterion, it can be checked
  from the leading principal minors alone, whereas positive semidefiniteness
  needs all principal minors.
\end{remark}

\subsection{Verification}
\label{sec:verification}

A certificate is finite data: the ordered facets of $\mathcal{C}$, $\beta$,
the entries of $Q$ and the coefficients of the $h_{i}$, all rational. The
verifier checks the hypotheses of Proposition~\ref{prop:certificate} for this
data in exact rational arithmetic \artref{verifier}:
\begin{enumerate}
  \item \emph{Identity.} $D$ is rebuilt from the stored facets by
    \eqref{eq:signed-volume}, both sides of \eqref{eq:certificate} are
    expanded, and every coefficient must agree.
  \item \emph{Positive definiteness.} By Sylvester's criterion, $Q$ is
    positive definite if and only if every leading principal minor is
    positive. The minors are computed exactly by the fraction-free
    elimination of Bareiss \cite{bareiss-1968}.
  \item \emph{Orientation and type.} Every ridge must lie in exactly two facets with
    opposite induced orientations, so the orientation is coherent. The facet
    vertex sets must also be exactly those of the type: for $C_{4}(7)$ the
    fourteen facets listed in Remark~\ref{rem:cyclic}, which the verifier
    rebuilds from Gale's condition, and for eight vertices those of the
    polytope of Remark~\ref{rem:eighttypes} \artref{verifier}. Without these
    checks the bound would still hold for $D$, but need not bound the volume
    of a polytope of the type.
\end{enumerate}
The stored certificate for $C_{4}(7)$ passes all three checks and has
$\beta = 48/5$ \artref{certificates}, \artref{testmethod}; so do the $36$
certificates for eight vertices, with $\beta = 78/5$
\artref{certificateseight}.

\subsection{How the certificates were found}

\begin{remark}
  The certificates were found numerically and then made exact. The proofs use
  none of the steps below, only the checks of
  Section~\ref{sec:verification} \artref{certscript}.
\end{remark}

\begin{remark}
  \label{rem:computation}
  The entries of $Q$ and the coefficients of the $h_{i}$ enter
  \eqref{eq:certificate} linearly, and the only other condition is that $Q$
  be positive semidefinite, so finding them is a semidefinite program. We
  solved it in floating point with Clarabel \cite{goulart-chen-2026} through
  CVXPY \cite{diamond-boyd-2016}, making the smallest eigenvalue of $Q$ as
  large as possible \artref{search}.
\end{remark}

\begin{remark}
  Changing the signs of two coordinates of every point, say $x_{i,a}$ and
  $x_{i,b}$ for all $i$, leaves $D$ unchanged. We grouped the monomials by how
  they change sign under these maps and took $Q$ block diagonal, with one
  block per group \artref{basis}. This made the program much smaller; the
  verifier checks each block.
\end{remark}

\begin{remark}
  We rounded every number to a rational with denominator $10^{9}$ and added
  the small leftover of \eqref{eq:certificate} to matching entries of $Q$, so
  that \eqref{eq:certificate} holds exactly. This is a simpler form of the
  rounding of Peyrl and Parrilo \cite{peyrl-parrilo-2008} \artref{rounding}.
  The verifier checks that $Q$ is still positive definite.
\end{remark}

\begin{remark}
  For $C_{4}(7)$, any $\beta < 4! \cdot \sqrt{3}/4 \approx 10.39$ with a
  certificate proves Corollary~\ref{cor:problem2}, and none exists below the
  maximum of $D$, which is at least the value of about $9.34$ found
  numerically \artref{ascent}, \artref{certificates}, \artref{testmethod}. We chose $\beta = 48/5$, giving the simple
  bound $2/5$, and did not look for the smallest possible $\beta$. For eight
  vertices, any common $\beta < 4! \cdot 2/3 = 16$ suffices for
  Corollary~\ref{cor:eight}. We chose $\beta = 78/5$, giving the bound
  $13/20$, above the largest numerical value of $D$ among the $36$ types,
  about $4! \cdot 0.601 \approx 14.4$ \artref{sweepeight}.
\end{remark}
